% Recall Studio project-status brief
% Place this file in the same Overleaf directory as the supplied PNG assets.
% Expected filenames are listed in the Product Walkthrough section below.

\documentclass[11pt,letterpaper]{article}

\usepackage[T1]{fontenc}
\usepackage[utf8]{inputenc}
\usepackage{lmodern}
\usepackage[margin=0.82in]{geometry}
\usepackage{microtype}
\usepackage[table]{xcolor}
\usepackage{graphicx}
\usepackage{booktabs}
\usepackage{tabularx}
\usepackage{array}
\usepackage{enumitem}
\usepackage{titlesec}
\usepackage{caption}
\usepackage{fancyhdr}
\usepackage{hyperref}
\usepackage{lastpage}

\definecolor{RecallInk}{HTML}{0B0E18}
\definecolor{RecallPanel}{HTML}{F3F4F8}
\definecolor{RecallLine}{HTML}{D8DCEF}
\definecolor{RecallViolet}{HTML}{7466D9}
\definecolor{RecallMint}{HTML}{438F7F}
\definecolor{RecallMuted}{HTML}{5E6472}
\definecolor{RecallWarm}{HTML}{A87026}
\definecolor{RecallDanger}{HTML}{B44254}

\hypersetup{
  colorlinks=true,
  linkcolor=RecallViolet,
  urlcolor=RecallMint,
  pdftitle={Recall Studio Project Status and Technical Brief},
  pdfauthor={Recall Studio}
}

\setlength{\parindent}{0pt}
\setlength{\parskip}{0.68em}
\setlist[itemize]{leftmargin=1.25em,itemsep=0.28em,topsep=0.25em}
\setlist[enumerate]{leftmargin=1.45em,itemsep=0.35em,topsep=0.25em}
\renewcommand{\arraystretch}{1.22}

\titleformat{\section}
  {\Large\bfseries\color{RecallInk}}
  {\thesection}{0.65em}{}
  [\vspace{0.2em}\color{RecallLine}\titlerule]
\titleformat{\subsection}
  {\large\bfseries\color{RecallInk}}
  {\thesubsection}{0.6em}{}
\titleformat{\subsubsection}
  {\normalsize\bfseries\color{RecallInk}}
  {\thesubsubsection}{0.5em}{}

\pagestyle{fancy}
\fancyhf{}
\fancyhead[L]{\small\sffamily\bfseries RECALL STUDIO}
\fancyhead[R]{\small\sffamily\color{RecallMuted} Project Status}
\fancyfoot[L]{\small\sffamily\color{RecallMuted} Pre-beta technical brief}
\fancyfoot[R]{\small\sffamily\color{RecallMuted} \thepage\ of \pageref{LastPage}}
\renewcommand{\headrulewidth}{0.4pt}
\renewcommand{\footrulewidth}{0pt}

\newcolumntype{Y}{>{\raggedright\arraybackslash}X}
\newcolumntype{P}[1]{>{\raggedright\arraybackslash}p{#1}}

\newcommand{\statuslabel}[2]{%
  \begingroup
  \setlength{\fboxsep}{4pt}%
  \colorbox{#1!12}{\textcolor{#1}{\sffamily\footnotesize\bfseries #2}}%
  \endgroup
}

\newcommand{\callout}[1]{%
  \begin{center}
    \setlength{\fboxsep}{10pt}%
    \colorbox{RecallPanel}{%
      \parbox{0.93\linewidth}{\color{RecallInk}#1}%
    }
  \end{center}
}

\newcommand{\missingasset}[1]{%
  \fcolorbox{RecallLine}{RecallPanel}{%
    \parbox[c][0.28\textheight][c]{0.88\linewidth}{%
      \centering\sffamily\color{RecallMuted}
      Image asset not found\\[0.5em]
      \texttt{\detokenize{#1}}%
    }%
  }%
}

\newcommand{\productimage}[2]{%
  \IfFileExists{#1}{%
    \includegraphics[width=\linewidth,height=0.68\textheight,keepaspectratio]{#1}%
  }{%
    \missingasset{#2}%
  }%
}

\newcommand{\productfigure}[3]{%
  \begin{figure}[p]
    \centering
    \productimage{#1}{#1}
    \caption{\textbf{#2} #3}
  \end{figure}
  \clearpage
}

\begin{document}

\begin{titlepage}
  \thispagestyle{empty}

  \IfFileExists{Logo.png}{%
    \includegraphics[width=25mm,height=25mm,keepaspectratio]{Logo.png}%
  }{%
    \fcolorbox{RecallViolet}{white}{%
      \parbox[c][22mm][c]{22mm}{\centering\sffamily\bfseries\color{RecallViolet}RECALL}%
    }%
  }

  \vspace{1.25in}

  {\fontsize{34}{39}\selectfont\bfseries\color{RecallInk} Recall Studio\par}
  \vspace{0.18in}
  {\Large\sffamily\color{RecallMuted} Project Status and Technical Brief\par}

  \vspace{0.55in}

  \callout{%
    \textbf{Recall Studio is local-first session memory for Ableton Live.}
    It records meaningful production activity, organizes it into a reviewable
    history, and helps a producer understand how a project changed over time.
  }

  \vfill

  \begin{tabular}{@{}ll@{}}
    \textbf{Product stage:} & Pre-beta development \\
    \textbf{Application version:} & 0.1.0 \\
    \textbf{Primary platform:} & Windows 10/11 \\
    \textbf{Status date:} & July 22, 2026 \\
  \end{tabular}

  \vspace{0.5in}

  {\small\color{RecallMuted}
  Prepared for developers, technical advisors, potential collaborators, and
  people evaluating the product.}
\end{titlepage}

\tableofcontents
\clearpage

\section{Executive Summary}

Recall Studio is a desktop companion for Ableton Live. It captures the work that
normally disappears between saved project files: devices added or removed,
parameter moves, clips, tempo changes, takes, and creative moments a producer
chooses to keep. The application turns that activity into a project history that
can be reviewed after the session.

The product is functional but not release-ready. The native application, local
database, project and take management, timeline, recap tools, release organizer,
notes, and diagnostic report flow are implemented. The capture layer currently
supports the bundled Max for Live bridge over UDP. A newer Ableton control-surface
integration over TCP is present in the current branch as an active prototype and
still needs sustained live validation before it should become the default path.

\callout{%
  \textbf{The core product idea is simple:} an Ableton project file remembers the
  result. Recall Studio is intended to remember the work that produced it.
}

The immediate challenge is no longer proving that events can reach a desktop app.
The challenge is earning trust over real, multi-hour sessions. That means validating
the new capture path, improving field diagnostics, removing silent failure modes,
and testing installation on machines that do not have development tools.

\section{The Product}

\subsection{The problem it addresses}

Music production is a long sequence of small decisions. A producer may find the
right bass tone through a short filter ride, delete a device that later becomes
important, or save several versions without remembering which one contained the
best idea. Ableton preserves the current project state, but it does not provide a
readable history of how that state developed.

Manual notes are useful when someone remembers to write them. In practice, the
best sessions are often the moments when note-taking stops. Cloud collaboration
tools solve a different problem and can be inappropriate for unreleased material.
Recall therefore keeps its working data on the producer's machine and captures in
the background.

\subsection{Product principles}

\begin{itemize}
  \item \textbf{Invisible during creation.} Capture must not interrupt the producer
        or make Ableton feel slower.
  \item \textbf{Local by default.} Session data is stored in a local SQLite database.
        Nothing is uploaded by the application.
  \item \textbf{The raw event log remains authoritative.} Timelines and summaries are
        derived views. Curation may hide noise, but it should not rewrite history.
  \item \textbf{Producer language over implementation language.} The interface talks
        about projects, takes, tracks, moves, and moments rather than packets or schemas.
  \item \textbf{Uncertainty is shown honestly.} When Ableton or a plug-in does not
        expose information, the product should identify the gap instead of inventing detail.
\end{itemize}

\section{Current Product State}

\subsection{Status at a glance}

\begin{center}
\small
\rowcolors{2}{RecallPanel}{white}
\begin{tabularx}{\linewidth}{P{0.19\linewidth} P{0.16\linewidth} Y Y}
  \toprule
  \textbf{Area} & \textbf{Status} & \textbf{What works now} & \textbf{Current boundary} \\
  \midrule
  Startup and Project Desk & \statuslabel{RecallMint}{Implemented} &
  Local producer profile, bridge state, project library, active take controls, folder connection, and project creation. &
  The guided first-run music-folder scan is not yet a complete onboarding flow. \\

  Projects and Versions & \statuslabel{RecallMint}{Implemented} &
  Projects, loose takes, scanned \texttt{.als} files, relinking, take status, and structural comparison between versions. &
  Parameter-level comparison and persistent cross-take track identity are planned. \\

  Timeline & \statuslabel{RecallMint}{Implemented} &
  Curated session activity, track lanes, parameter moves, creative moments, activity blocks, filtering, and session recap data. &
  Arrangement-position anchoring and the planned Sound Story view are not implemented. \\

  Recap and export & \statuslabel{RecallMint}{Implemented} &
  Session summary plus Markdown, text, JSON, and printable export paths in the sharing pipeline. &
  Export behavior still needs clean-machine and real-project validation. \\

  Project Organizer & \statuslabel{RecallMint}{Implemented} &
  Album, EP, and single organization; track ordering; source sets; bounce versions; waveform analysis; loudness and peak measurements; timed comments; final-version selection; and private preview export. &
  Its exact role in the first beta scope should be decided so it does not distract from capture reliability. \\

  Notes and Reference & \statuslabel{RecallWarm}{Partial} &
  Autosaving producer notes and a built-in production reference are available. &
  Notes still use browser local storage and must move to SQLite before beta. \\

  Diagnostics & \statuslabel{RecallWarm}{Partial} &
  A Report a Problem dialog can gather connection details, bridge metrics, recent frontend errors, and a producer note for copy or save. &
  Persistent file logging, a dedicated diagnostics view, and a React error boundary remain open. \\

  Capture transport & \statuslabel{RecallWarm}{Transition} &
  Max for Live over UDP remains available. A Python control-surface path over TCP now feeds the same Rust pipeline. &
  The TCP path is still labeled as a spike in source and needs marathon-session and compatibility testing. \\
  \bottomrule
\end{tabularx}
\end{center}

\subsection{What has been proven}

The application can receive structured events, normalize multiple protocol shapes,
classify event importance, write events to SQLite, rebuild typed project data, and
render that history in React. Folder scanning can discover existing \texttt{.als}
files and create takes without modifying the Ableton projects themselves.

The ingestion pipeline includes bounded queues, transactional persistence, sequence
gap detection, panic isolation per event, and an enlarged UDP receive buffer. The
project's stress work showed why this matters: packet loss occurred before the
application queue when Windows exhausted its default socket buffer. Increasing the
receive buffer and moving work off the receive loop corrected that measured failure
mode in the test harness.

This is meaningful engineering progress, but it is not the same as a completed beta.
The remaining standard is operational: Recall must run beside Ableton for hours,
capture the right details, and explain any failure clearly enough to diagnose later.

\subsection{Known limitations}

\begin{itemize}
  \item The product is Windows-first. macOS packaging, signing, and notarization are
        separate work.
  \item Some VST plug-ins do not expose preset names or complete state through Ableton's
        public object model. Recall can record only what Live makes visible.
  \item The control-surface prototype was written against Ableton Live 12's embedded
        Python environment. Live 11 compatibility has not been established for that path.
  \item The application does not yet have production file logging, signed installers,
        an updater, or continuous integration.
  \item The Tauri content security policy is currently unset, and package metadata
        remains at development defaults in several places.
  \item Several shared Rust states still use poisoning mutexes with \texttt{expect} calls.
        A failure while a lock is held can make later commands fail until restart.
  \item Sound Story, arrangement anchoring, first-run scanning, and parameter-level
        take differences are specified or planned, not shipped features.
\end{itemize}

\section{Product Walkthrough}

The screenshots in this section show the current interface. Each image is presented
as product evidence, not as a claim that every workflow has completed release testing.

\productfigure{Landing Page.png}{Startup and local profile.}{The opening screen shows local identity, library readiness, and Ableton bridge state before entering the Project Desk.}

\productfigure{Nav Bar.png}{Application navigation.}{Projects, Recap, Timeline, Organizer, Notes, and Reference are persistent surfaces. The top-left Recall mark returns to startup.}

\productfigure{Project Manager.png}{Project Desk.}{The project library groups takes by song, supports connected folders, exposes active capture state, and provides direct access to recap and timeline views.}

\productfigure{Timeline_Capture.png}{Session capture timeline.}{Captured events are organized into a readable working history rather than displayed as an unfiltered transport log.}

\productfigure{Track_Changes_Timeline.png}{Track-level changes.}{The timeline groups activity by musical context so parameter moves and structural changes can be read after the session.}

\productfigure{project Organizer.png}{Project Organizer.}{A release can contain multiple songs, source Ableton sets, bounce versions, measurements, comments, and a selected final version.}

\productfigure{Unfilled Project Organizer.png}{Organizer empty state.}{The incomplete state makes missing source sets, mixes, or release information explicit rather than presenting an apparently finished release.}

\productfigure{Album Preview.png}{Private release preview.}{Final bounce selections can be assembled into a portable listener page with cover art, playback, track order, and comments.}

\productfigure{Notes.png}{Producer Notes.}{Notes live outside an individual session and autosave locally. Native SQLite persistence is still required before beta.}

\productfigure{Reference.png}{Production Reference.}{A built-in reference keeps common production terminology and practical guidance available without leaving the application.}

\section{Technical Breakdown}

\subsection{Architecture at a glance}

\begin{center}
\small
\begin{tabularx}{\linewidth}{P{0.20\linewidth} P{0.23\linewidth} Y}
  \toprule
  \textbf{Layer} & \textbf{Technology} & \textbf{Responsibility and rationale} \\
  \midrule
  Ableton integration & Max for Live JavaScript bridge; Python MIDI Remote Script prototype &
  Observes Live's object model and emits changes. Max for Live is the established path. The control-surface prototype explores broader capture with lower traversal cost and without UDP snapshot-size limits. \\

  Local transport & UDP on \texttt{127.0.0.1:9000}; TCP on \texttt{127.0.0.1:9001} &
  UDP keeps the existing bridge simple and low overhead. TCP carries newline-delimited JSON from the prototype and supports whole-set snapshots larger than one UDP datagram. Both remain local to the machine. \\

  Native application & Tauri 2 and Rust &
  Owns process boundaries, filesystem access, dialogs, transport listeners, persistence, and commands exposed to the frontend. Tauri keeps the desktop shell small while allowing native Rust control over reliability-sensitive work. \\

  Persistence & SQLite through \texttt{rusqlite} &
  Stores the append-only event history, projects, sessions, creative moments, organizer metadata, and migrations in a local database. SQLite fits the single-user, local-first model and supports transactional writes without a separate service. \\

  Projection layer & Rust schema projection &
  Rebuilds project, track, device, parameter, and change views from stored events. Derived views can evolve without discarding the original capture record. \\

  Interface & React 19, TypeScript, and Vite &
  Renders project management, versions, timeline, recap, organizer, notes, reference, and diagnostics. Tauri IPC keeps native operations behind explicit commands. \\
  \bottomrule
\end{tabularx}
\end{center}

\subsection{Event flow}

\begin{enumerate}
  \item Ableton exposes tracks, devices, parameters, clips, tempo, and project context
        through its object model.
  \item The active bridge converts relevant observations into Recall Protocol JSON.
        The established Max for Live path sends UDP datagrams. The newer control-surface
        path queues newline-delimited JSON to a background TCP sender so socket work does
        not block Live's main thread.
  \item Rust validates the protocol envelope and normalizes older or nested field shapes
        into one canonical event representation.
  \item Events are assigned a priority. Creative and structural changes are protected.
        High-frequency telemetry may be shed under overload.
  \item A bounded 4,096-event channel separates receiving from database work. The worker
        persists as many as 256 events per transaction.
  \item SQLite retains the raw record. Projection code derives typed entities, parameter
        histories, summaries, versions, and timeline material.
  \item React requests those views through Tauri commands. Export occurs only when the
        producer explicitly chooses to create a recap or release preview.
\end{enumerate}

\subsection{Reliability controls currently in code}

\begin{itemize}
  \item The UDP socket requests an 8 MiB receive buffer to absorb project-load bursts.
  \item Sequence tracking detects gaps independently for each sending device.
  \item The persistence channel is bounded at 4,096 events and never blocks the receive loop.
  \item The in-memory live-event buffer is capped at 50,000 entries; SQLite remains authoritative.
  \item Each UDP packet and TCP event is processed behind panic isolation so one malformed
        payload does not end the capture thread.
  \item The control-surface sender has a bounded 2,048-event queue and discards the oldest
        backlog first when the application is unavailable.
  \item Parameter listening in the prototype is limited to the focused device and at most
        128 parameters. A 350 ms settle window turns a dense knob ride into one meaningful move.
\end{itemize}

\subsection{Data and privacy posture}

Recall does not require an account and does not upload session data. Transport is
bound to loopback addresses. Core history and organizer metadata live in SQLite;
larger waveform and cover assets are stored as managed files in the application's
data directory. Notes are the current exception because they still use local storage.

The application is read-only with respect to Ableton's musical project. It observes
and records activity but does not inject locators, rewrite project history, or modify
the \texttt{.als} file. Export is an explicit user action.

\subsection{Verified build and test state}

The following results were verified against the current working tree on July 22, 2026:

\begin{itemize}
  \item TypeScript compilation and the Vite production build complete successfully.
  \item 112 frontend tests pass across 9 Vitest files.
  \item 55 Rust tests pass, including storage migration, schema projection, transport
        normalization, sequence-gap detection, organizer persistence, and preview export.
  \item The repository does not yet run these checks in continuous integration, so the
        verified local result is stronger than the current release process around it.
\end{itemize}

\section{What's Next}

\subsection{Immediate engineering priorities}

\begin{enumerate}
  \item \textbf{Choose and validate the beta capture path.} Run the TCP control-surface
        prototype beside real Ableton projects for multi-hour sessions. Confirm CPU impact,
        reconnection behavior, listener cleanup, snapshot completeness, and Live 11/12 scope.
        Then decide whether TCP replaces Max for Live or ships as a separate capture tier.
  \item \textbf{Make failures observable.} Add persistent native file logging, expose the
        full bridge metrics in a dedicated diagnostic view, and retain the existing report
        generator as the handoff format for bug reports.
  \item \textbf{Remove silent failure modes.} Replace poisoning mutex usage, add a React
        error boundary around application surfaces, and provide a clear startup error when
        database initialization or a listener bind fails.
  \item \textbf{Dogfood the core loop.} Use Recall during every real session for at least
        two weeks. Track data loss, Ableton responsiveness, wrong project assignment, and
        whether the timeline actually answers a question the producer had forgotten.
  \item \textbf{Complete first-run value.} Turn the existing folder scanner into a guided
        first-run flow so a producer can see useful project history before a new capture has
        accumulated.
\end{enumerate}

\subsection{Before a friend beta}

\begin{itemize}
  \item Build the planned Sound Story view from existing event data and validate it on real tracks.
  \item Move Notes from local storage into SQLite with an additive migration.
  \item Add continuous integration for the Rust tests, frontend tests, type checking, and build.
  \item Synchronize application versions and replace development package metadata.
  \item Test the NSIS installer and bridge setup on a clean Windows machine with no developer tools.
  \item Set a real Tauri content security policy and document known plug-in visibility limits.
\end{itemize}

\subsection{After the first stable release}

Arrangement-position anchoring is the strongest planned extension. It would place
changes on the song's bars-and-beats ruler instead of only on session time. That work
requires a protocol update and its own Live performance validation, so it should not
be pulled into the initial release path.

Other later candidates include parameter-level take differences, persistent identity
for tracks across renamed versions, richer collaboration exports, a signed macOS build,
and an updater. AI-generated summaries remain a possible later layer over trustworthy
captured data, not a substitute for getting capture and history correct.

\section{Questions or Claims to Verify}

\begin{enumerate}
  \item Is the Python control-surface integration intended to replace the Max for Live
        bridge, or should both remain supported capture tiers?
  \item Does the control-surface path support Ableton Live 11, or is Live 12 the accurate
        requirement for the first version that uses it?
  \item Has the TCP snapshot and focused-device listener path completed a real four-hour
        session with instruments active and no measurable effect on Live?
  \item Is the Project Organizer part of the first friend beta, or should the beta remain
        focused on project history, versions, timeline, recap, and Sound Story?
  \item Which export formats have been visually reviewed on a clean packaged build, as
        opposed to being implemented and covered by unit tests?
  \item The README and several planning documents still describe UDP and Max for Live as
        the only active capture path. They should be updated once the transport decision is final.
\end{enumerate}

\end{document}
